\chapter{Introduction}
\label{ch:intro}

\section{Motivation}
The electrification of traction, industrial actuation and household appliances
has shifted a large fraction of product validation from the vehicle or the
appliance itself onto the dynamometer bench. A dyno can impose a controlled
mechanical load, record shaft torque and speed, and---when coupled to an
inverter---exercise the full current, flux and speed loops of a
field-oriented controller~\cite{isermann2014engine,krishnan2010pmsm}. The
value of that facility, however, is only as high as the quality of the
digital interface that binds the bench software to the unit under test
(UUT).

In the laboratory that hosted this work, the UUT is an STMicroelectronics
Stellar-E microcontroller of the SR5E1 family~\cite{ststellar}, running the
Motor Control software development kit (MCSDK)~\cite{stmcsdk}. The physical
link is a CAN (and CAN FD) bus~\cite{iso11898,bosch1991can}. Historically,
bench commands were translated into the internal ST Motor Control protocol
by a firmware layer named DYNO2DW---``Dyno to Dishwasher''---a name that
betrays the appliance origin of the protocol profile. The layer was
functional enough for manual bring-up, but it was not a contract that a
test automation engineer could trust: identifiers collided, a single
\texttt{EXECUTE\_COMMAND} frame could request both start and stop, PI gains
could not be written atomically, $I_d$ and $I_q$ could not be ramped
independently, and the DBC file that Vector CANalyzer used to decode the bus
had drifted away from the C implementation.

The practical consequence is easy to state. An operator who wants to sweep
speed from $0$ to $n^\star$ in $T$ milliseconds, hold a torque-producing
current $I_q$, detune the speed PI, and log heatsink temperature, should
press four well-labelled controls---not compose eight hexadecimal frames
while watching a trace window. Closing that gap is the subject of this
thesis.

\section{Problem statement}
Three artefacts must stay consistent if a CAN-based dyno interface is to be
automatable:

\begin{enumerate}[leftmargin=1.6em]
  \item the \emph{firmware} that unpacks incoming frames, calls the Motor
        Control API, and emits diagnostics;
  \item the \emph{DBC} that names every identifier, places every signal and
        documents every scale factor;
  \item the \emph{operator surface}---here Vector CANalyzer plus a Python
        GUI---that hides those identifiers behind physical units.
\end{enumerate}

At the start of the project none of the three was in that state. The SR5E1
CAN library implemented a subset of the Motor Control protocol, with
offset-based identifier mapping ($+$\texttt{0x30} for commands,
$+$\texttt{0x40} for diagnostics) that had never been fully reconciled with
the DBC. Reception lived in an interrupt service routine and was therefore
difficult to inspect without a hardware debugger. Transmission of the
periodic status stack was coupled to the Frame Communication Protocol (FCP)
state machine in a way that silently dropped updates when the transmitter
was busy. Several protocol codes required by the bench---independent $I_d$
and $I_q$ references, a flux ramp, packed PI writes---were either missing or
wired to the wrong identifier.

\section{Objectives}
The work was organised around three objectives, corresponding to the three
artefacts above.

\paragraph{O1 --- Validate and correct the SR5E1 CAN library.}
Exercise every incoming command and every outgoing diagnostic on hardware.
Prove, with TRACE32 breakpoints and PCAN-View traces, that each CAN
identifier maps onto the intended Motor Control call. Remove contradictory
control paths and colliding response identifiers.

\paragraph{O2 --- Rebuild the DBC so that it is a faithful dictionary of the
bus.}
Add the new messages required by O1, enforce the $+$\texttt{0x30}/$+$\texttt{0x40}
offset convention, and correct the structural error that treated
acknowledgement codes as CAN identifiers rather than payload values.

\paragraph{O3 --- Provide an operator GUI on top of Vector CANalyzer.}
Use CAPL and the CANalyzer COM API so that a Python interface can start and
stop the motor, program speed and current ramps, tune PI gains and plot
telemetry, without exposing hexadecimal layouts to the user.

\section{Contributions}
The concrete contributions of this thesis are the following.

\begin{itemize}[leftmargin=1.6em]
  \item A documented architecture of the DYNO2DW translation layer, including
        the offset rule, the ten-message diagnostic stack and the
        transmission/reception call graph
        (\cref{ch:arch}).
  \item A mutually exclusive rewrite of \texttt{EXECUTE\_COMMAND} that
        cannot emit both \texttt{START\_MOTOR} and \texttt{STOP\_MOTOR} for
        the same frame, together with a clean handling of reset, fault
        acknowledgement and encoder alignment (\cref{ch:dev}).
  \item Independent $I_d$ and $I_q$ set-points
        (\texttt{MC\_PROTOCOL\_CODE\_SET\_CURRENT\_ID\_REF},
        \texttt{MC\_PROTOCOL\_CODE\_SET\_CURRENT\_IQ\_REF}) and a newly
        implemented flux ramp (\texttt{UI\_ExecFluxRamp} /
        \texttt{MCI\_ExecFluxRamp} / \texttt{STC\_ExecCurrentRamps}) that
        mirrors the existing torque ramp (\cref{ch:dev}).
  \item A 32-bit packing convention for PI registers, so that a single
        \texttt{SET\_REG} writes both $K_p$ (or $K_i$) and its
        power-of-two divisor (\cref{ch:dev}).
  \item Separation of ACK/NACK (identifier \texttt{0x10}, payload
        \texttt{0xF0}/\texttt{0xFF}) from the \texttt{GET\_REG} response
        (identifier \texttt{0x11}), implemented in
        \texttt{CFCP\_TX\_IRQ\_Handler} by dispatching on frame size
        (\cref{ch:dev}).
  \item A case study of a specification-versus-implementation discrepancy:
        in the SR5E1 library, byte~0 of a \texttt{SET\_REG} frame is the
        register identifier, not a length field as the dishwasher protocol
        document suggests (\cref{ch:val}).
  \item A DBC aligned on the firmware, with a corrected acknowledgement
        message and a reserved identifier window for new commands
        (\cref{ch:dbc}).
  \item A Python GUI, driven through \texttt{win32com.client}, covering
        motor start/stop, speed ramps, PI panels, $I_d$/$I_q$ control and
        live charts, with CAPL-side ready-flags that prevent incomplete
        multi-parameter frames from reaching the bus (\cref{ch:gui}).
\end{itemize}

\section{Methodology}
The work followed a closed-loop debug methodology rather than a
spec-first rewrite. Each defect was first reproduced on the bench by
injecting a CAN frame with PCAN-View over a PCAN-USB FD adapter. The
corresponding FDCAN interrupt was then trapped in Lauterbach TRACE32, so
that payload unpacking, the DYNO2DW switch statement, the Motor Control
call and the TX response could be observed on a single stepping path.
Only after a behaviour was understood in that environment was the C
source---or the DBC, or the CAPL---changed. The GUI was introduced last,
once the bus contract was stable, because an operator surface that sits on
a moving DBC is more expensive to maintain than a hexadecimal injector.

This ordering is deliberate. In embedded protocol work, the firmware is the
source of truth: a DBC that disagrees with the MCU will decode the bus
wrongly, and a GUI that disagrees with the DBC will send frames the MCU
silently nacks. Keeping that hierarchy explicit is, in itself, one of the
engineering results of the project.

\section{Structure of the thesis}
\Cref{ch:bg} reviews the technical background that a reader needs in order
to follow the rest of the text: CAN and CAN FD, DBC files, field-oriented
control of PMSM machines, the Stellar-E platform and the ST Motor Control
SDK, and the role of a dynamometer bench. \Cref{ch:arch} describes the
legacy SR5E1 CAN library as it was found. \Cref{ch:dev} presents the
firmware developments. \Cref{ch:val} reports the hardware validation
campaign and the protocol-discrepancy case study. \Cref{ch:dbc} documents
the DBC rebuild. \Cref{ch:gui} describes the CANalyzer--Python integration
and the operator panels. \Cref{ch:conc} summarises the results and outlines
future work. An appendix collects the longer firmware and CAPL listings.
